\section{Related Work}
\label{sec:related}

We organize prior work around the paper's three claims: that derived state outlives authorization
changes, that response metadata reveals hidden documents, and that provenance-based read-time
enforcement closes the first without opening the second. \Cref{tab:related} positions the closest
work on the dimensions that matter here.

\begin{table*}[t]
\centering
\caption{Closest prior work. Covers: which derived artifacts the work addresses. Per-principal:
removal for one user while others keep access (revocation), as opposed to removal for everyone
(deletion). Generated: handles artifacts written by a generative model, which carry no source
identifiers. Lifetimes: measures how long a stale copy stays servable.}
\label{tab:related}
\footnotesize
\setlength{\tabcolsep}{3pt}
\resizebox{\textwidth}{!}{%
\begin{tabular}{@{}lcccccccl@{}}
\toprule
Work & Index & Cache & Memory & Revocation & Per-principal & Generated & Lifetimes & Kind \\
\midrule
Zanzibar~\cite{pang2019zanzibar} & -- & authz cache & -- & yes & yes & -- & -- & mechanism \\
Permission-aware RAG~\cite{jeong2025permissionrag}, HoneyBee~\cite{zhong2026honeybee} & yes & -- & -- & partly & yes & -- & -- & mechanism \\
Thoth~\cite{elnikety2016thoth}, Qapla~\cite{mehta2017qapla} & yes & -- & -- & -- & yes & -- & -- & mechanism \\
DELF~\cite{cohngordon2020delf}, K9db~\cite{albab2023k9db} & stores & stores & -- & -- & -- & -- & -- & mechanism \\
Tavakoli and Sanderson~\cite{tavakoli2026deletion} & yes & -- & -- & -- & -- & -- & -- & measurement \\
FIDES~\cite{costa2026fides}, CaMeL~\cite{debenedetti2025camel} & -- & -- & agent state & -- & -- & yes & -- & mechanism \\
MEXTRA~\cite{wang2025mextra} & -- & -- & yes & -- & -- & yes & -- & attack \\
SAAM~\cite{crampton2006saam}, recycling~\cite{wei2009recycling} & -- & authz cache & -- & policy change & yes & -- & -- & mechanism \\
Liu et al.~\cite{liu2015revocation} (web PKI) & -- & CRL/OCSP & -- & yes & -- & -- & yes & measurement \\
\textbf{This work} & yes & yes & yes & yes & yes & yes & yes & measurement + comparison \\
\bottomrule
\end{tabular}}
\end{table*}

\subsection{Derived state outlives authorization changes}
\paragraph{Access control for search and retrieval.}
Document-level security in enterprise search predates RAG. Bailey et al.\ compare copying
permissions into the index with checking each result at query time~\cite{bailey2006secure};
fine-grained access control by query rewriting is the database analogue~\cite{rizvi2004query}. RAG
systems filter retrieval with permissions derived from an IAM system~\cite{jeong2025permissionrag},
for sensitive patient data~\cite{chen2025integrating}, or with dynamically partitioned vector
indexes for role-based filtering~\cite{zhong2026honeybee}; products expose the same pattern as
metadata filters~\cite{doc_azure_trimming,doc_qdrant,doc_pinecone,doc_chroma}. Filtered
nearest-neighbor search treats the pre- versus post-filter choice as an efficiency
problem~\cite{gollapudi2023filtered,patel2024acorn}; we show it is also a security choice. A
materialized ACL copy is a materialized view, stale for as long as its maintenance
lags~\cite{gupta1995views}. None of these designs addresses artifacts derived before a permission
changed.

\paragraph{Revocation and stale authorization.}
Complete mediation requires every access to be checked against current
authority~\cite{saltzer1975protection}. Revocation is an old problem for capabilities, where
selective revocation needs indirection~\cite{redell1974selective}, and for cryptographic file
systems, where \emph{lazy revocation} defers re-encryption until the next
write~\cite{kallahalla2003plutus,backes2005lazy}; our read-time check is the opposite trade-off,
paying on every read so that no stale copy is served. Secondary and approximate authorization
models cache and infer decisions for availability and must bound how stale they
are~\cite{crampton2006saam,wei2009recycling}. Zanzibar binds checks to consistency tokens to avoid
the ``new enemy'' problem~\cite{pang2019zanzibar}; our artifacts are derived \emph{before} the
revocation and re-served after it with no intervening write. Measurement work on certificate
revocation~\cite{liu2015revocation} is the methodological template for measuring how long a revoked
credential stays usable; we apply that view to RAG artifacts.

\paragraph{Deletion, unlearning and derived data.}
A system must forget data together with everything computed from it~\cite{cao2015towards}, formally
captured by deletion-compliance definitions~\cite{garg2020formalizing}. DELF makes cascading deletion
first-class at Facebook and reports that developers routinely omit derived-data deletions: before
enforcement they handled deletion for only 20\% of new edge types~\cite{cohngordon2020delf}. K9db and
Edna build compliant deletion and disguising into web-application storage~\cite{albab2023k9db,tsai2023edna},
and GDPR compliance has measurable cost in databases~\cite{shastri2020gdpr}. Unlearning removes data
from trained models~\cite{bourtoule2021machine} and, for RAG, from the knowledge
base~\cite{wang2024unlearningrag}; Tavakoli and Sanderson audit whether RAG systems forget deleted
material~\cite{tavakoli2026deletion}. Revocation differs from deletion in the way that matters for
derived state: the document stays available to other principals, so a derived copy must be
filtered per principal on every read rather than removed.

\paragraph{Chat memory.}
LLM applications keep conversation state as recent turns plus a rolling summary written by the
model~\cite{wang2023recursively,doc_langchain_summary,doc_llamaindex_memory}, or as managed long-term
memory~\cite{packer2023memgpt,park2023generative,zhong2024memorybank}. Memory can be extracted by a
crafted query~\cite{wang2025mextra}. No prior work we found asks what memory serves after the user's
permissions change, which is where our measurements find the longest-lived leaks.

\subsection{Metadata reveals hidden documents}
Statistics computed over documents a user cannot read leak information about them: under a
system-wide index with post-filtering, an unprivileged user recovers collection statistics from
relevance scores~\cite{buttcher2005security}. Cross-site search attacks infer the existence of private
records from search timing~\cite{gelernter2015xs}, and web applications leak state through response
size and timing~\cite{chen2010sidechannel,bortz2007exposing,felten2000timing}. In LLM serving, prompt
and KV caches create timing channels~\cite{gu2025auditing,song2025earlybird,wu2025know,zheng2024inputsnatch},
speculative and efficient inference leaks through timing~\cite{carlini2024remote}, token lengths leak
encrypted responses~\cite{weiss2024what}, and semantic caches are vulnerable to key
collisions~\cite{zhang2026keycollision,afiffy2026semanticcache} and motivate user-scoped
designs~\cite{gill2024meancache}. Membership inference against the RAG datastore uses the generated
answer~\cite{anderson2025my,li2024generating,naseh2025riddle,liu2025mask}, and extraction attacks
recover datastore content~\cite{zeng2024good,qi2025follow,dimaio2024pirates}; embeddings themselves
can be inverted~\cite{morris2023embeddings}. Defenses include differentially private
retrieval~\cite{koga2024dprag} and synthetic datastores~\cite{zeng2024sage}. Those attacks target a
datastore the model can retrieve from. Our attacker cannot retrieve the target at all and asks
whether metadata or shared state betrays it; we report insiders and outsiders separately because
shared caches give the two very different power, and we use the low-false-positive operating point
of Carlini et al.~\cite{carlini2022mia}.

\subsection{Provenance-based read-time enforcement}
Provenance-tainted revocation is dynamic information-flow tracking~\cite{myers1997decentralized,sabelfeld2003language};
operating systems~\cite{zeldovich2006making,krohn2007information} and applications~\cite{yip2009improving,enck2010taintdroid}
attach labels to data and propagate them through computation, and provenance-based access control
names the model in which decisions depend on derivation history~\cite{park2012pbac,buneman2001why}.
Thoth enforces declarative policies on derived data inside a Lucene-based retrieval system, with a
per-task taint set~\cite{elnikety2016thoth}; Qapla enforces policies on database queries without
trusting the application~\cite{mehta2017qapla}. For LLM systems, information-flow control now spans
agent planners~\cite{costa2026fides,debenedetti2025camel,wu2024fsecure,zhong2025rtbas}, isolation
architectures~\cite{wu2025isolategpt}, contextual-integrity agents~\cite{bagdasarian2024airgap},
access control over model parameters~\cite{tiwari2024ifc,jayaraman2025permissioned}, and permissive
analyses that reduce over-tainting by asking which inputs actually influenced an
output~\cite{siddiqui2024permissive}; context attribution~\cite{cohenwang2024contextcite} is the
corresponding tool. These works establish the mechanism. Two things differ here. Our unit of taint
is the persistent artifact (a cache entry, a memory turn, a summary sentence) rather than a task or
a plan, and our question is what the mechanism costs compared with coarser invalidation, measured
on in-session utility. The usual limit of information-flow control, that a principal who read data
can re-type it, also bounds our defense, and we measure it.

\paragraph{Search protocol.}
On 14 October 2026 we ran 21 targeted queries through general web search, the OpenAlex and Crossref
APIs, and publisher and proceedings pages (ACM DL, IEEE Xplore, USENIX, Springer, PMLR, NeurIPS). The
queries covered permission-aware and multi-tenant RAG, RAG unlearning and deletion, stale authorization
and revocation in distributed systems, information-flow control for LLM agents, chat-memory privacy,
semantic-cache and KV-cache side channels, relevance-score leakage in multi-user search, backward-private
searchable encryption and revocable encryption, and a direct sweep for post-revocation leakage in RAG.
The search was targeted rather than systematic; work on agent memory, authorization caching, filtered
vector search and semantic caches was added afterwards from the authors' knowledge of these areas. We
found no peer-reviewed work that measures post-revocation leakage per derived artifact in a RAG pipeline.
